\paragraph{Poisson--Cauchy 粗粒化二阶正规形下的剖面极限与散度常数闭包}
以下结论统一沿用命题 \ref{con:xi-poisson-coarsegraining-second-order-normal-form} 的记号。设
\[
\tilde h_t(x):=(\mathsf P_t*\tilde\nu)(x),\qquad
\tilde g_t(x):=\mathsf P_t(x-\bar\gamma),\qquad
\mathsf P_t(x):=\frac{1}{\pi}\frac{t}{x^2+t^2},
\]
其中 \(\tilde\nu\) 为 \(\RR\) 上概率测度，\(\bar\gamma=\int \gamma\,\dd\tilde\nu(\gamma)\)，\(\sigma^2=\int(\gamma-\bar\gamma)^2\,\dd\tilde\nu(\gamma)\in(0,\infty)\)。
并记
\[
r_t(x):=\frac{1}{\pi}\frac{t\bigl(3(x-\bar\gamma)^2-t^2\bigr)}{\bigl((x-\bar\gamma)^2+t^2\bigr)^3}
=\frac12\,\partial_x^2\mathsf P_t(x-\bar\gamma),
\]
使得
\[
\tilde h_t=\tilde g_t+\sigma^2 r_t+\varepsilon_t,
\]
且在三阶中心绝对矩有限时有 \(\|\varepsilon_t\|_\infty=O(t^{-4})\)、\(\|\varepsilon_t\|_{L^1}=O(t^{-3})\)。

\begin{conclusion}[\texorpdfstring{$L^p$}{Lp} 族锐渐近与普适常数谱]\label{con:xi-poisson-lp-spectrum-sharp-family}
若
\[
\int_{\RR}|\gamma-\bar\gamma|^3\,\dd\tilde\nu(\gamma)<\infty,
\]
则对任意 \(p\in(1,\infty)\) 有
\[
\lim_{t\to\infty}
\frac{t^{3-\frac1p}}{\sigma^2}\,
\|\tilde h_t-\tilde g_t\|_{L^p(\RR)}
=
\frac1\pi\left(\int_{\RR}\frac{|3y^2-1|^p}{(1+y^2)^{3p}}\,\dd y\right)^{\!1/p},
\]
并且 \(p=\infty\) 端点满足
\[
\lim_{t\to\infty}
\frac{t^3}{\sigma^2}\,
\|\tilde h_t-\tilde g_t\|_{L^\infty(\RR)}
=\frac1\pi.
\]
特别地，
\[
\lim_{t\to\infty}
\frac{t^{5/2}}{\sigma^2}\,
\|\tilde h_t-\tilde g_t\|_{L^2(\RR)}
=\frac{\sqrt3}{4\sqrt\pi}.
\]
\end{conclusion}
\begin{proof}[证明]
由
\(
\tilde h_t-\tilde g_t=\sigma^2 r_t+\varepsilon_t
\)
与余项估计可得 \(\|\varepsilon_t\|_{L^p}=o(\|r_t\|_{L^p})\)（\(p\in(1,\infty]\)）。
令 \(y=(x-\bar\gamma)/t\)，则
\[
r_t(\bar\gamma+ty)=\frac{1}{\pi t^3}\frac{3y^2-1}{(1+y^2)^3},
\]
从而
\[
\|r_t\|_{L^p(\RR)}
=t^{-3+1/p}\frac1\pi\left(\int_{\RR}\frac{|3y^2-1|^p}{(1+y^2)^{3p}}\,\dd y\right)^{1/p}\quad(p<\infty),
\]
\[
\|r_t\|_{L^\infty(\RR)}
=\frac{1}{\pi t^3}\sup_{y\in\RR}\frac{|3y^2-1|}{(1+y^2)^3}
=\frac{1}{\pi t^3}.
\]
结合扰动项即可得到极限式。再由
\(
\int_{\RR}\frac{(3y^2-1)^2}{(1+y^2)^6}\dd y=\frac{3\pi}{16}
\)
推出 \(p=2\) 常数。证毕。
\end{proof}

\begin{conclusion}[自相似剖面极限与强收敛]\label{con:xi-poisson-selfsimilar-profile-strong-limit}
在结论 \ref{con:xi-poisson-lp-spectrum-sharp-family} 的条件下，定义
\[
\Delta_t(y):=\frac{t^3}{\sigma^2}
\Bigl(\tilde h_t(\bar\gamma+ty)-\tilde g_t(\bar\gamma+ty)\Bigr),\qquad y\in\RR.
\]
则
\[
\Delta_t(y)\xrightarrow[t\to\infty]{}\Delta_\infty(y)
:=\frac1\pi\frac{3y^2-1}{(1+y^2)^3}
\]
在 \(\RR\) 上一致成立；并且对任意 \(p\in[1,\infty)\) 有
\[
\Delta_t\to\Delta_\infty\quad\text{于 }L^p(\RR,\dd y).
\]
\end{conclusion}
\begin{proof}[证明]
由二阶正规形，
\[
\Delta_t(y)=t^3r_t(\bar\gamma+ty)+\frac{t^3}{\sigma^2}\varepsilon_t(\bar\gamma+ty)
=\Delta_\infty(y)+\frac{t^3}{\sigma^2}\varepsilon_t(\bar\gamma+ty).
\]
而 \(\|\varepsilon_t\|_\infty=O(t^{-4})\) 给出
\[
\sup_{y\in\RR}\left|\frac{t^3}{\sigma^2}\varepsilon_t(\bar\gamma+ty)\right|=O(t^{-1})\to0,
\]
故一致收敛成立。
又 \(\Delta_\infty\in L^p(\RR)\)（任意有限 \(p\)），并且
\(
|\Delta_t-\Delta_\infty|\le Ct^{-1}
\)
一致有界，支配收敛推出 \(L^p\) 收敛。证毕。
\end{proof}

\begin{conclusion}[\texorpdfstring{$D_{\mathrm{KL}}$}{KL} 的两项锐展开与显式六阶系数]\label{con:xi-poisson-kl-two-term-sharp-explicit}
若进一步满足
\[
\int_{\RR}|\gamma-\bar\gamma|^5\,\dd\tilde\nu(\gamma)<\infty,
\]
并记中心矩
\(
\mu_k:=\int(\gamma-\bar\gamma)^k\,\dd\tilde\nu(\gamma)\ (k\ge3),
\)
则
\[
D_{\mathrm{KL}}(\tilde h_t\mid \tilde g_t)
=\frac{\sigma^4}{8t^4}
+\frac{1}{t^6}\left(
\frac{\sigma^6}{64}-\frac{\sigma^2\mu_4}{8}+\frac{3\mu_3^2}{32}
\right)
+o(t^{-6}).
\]
\end{conclusion}
\begin{proof}[证明]
该式即定理 \ref{thm:xi-cauchy-poisson-kl-sharp-asymptotics}（等价写法见
\ref{thm:xi-poisson-kl-sixth-order-refined-binomials}）
的同参数化表达。证毕。
\end{proof}

\begin{conclusion}[Rényi 相对熵主系数的全参数普适性]\label{con:xi-poisson-renyi-alpha-over-8}
在结论 \ref{con:xi-poisson-lp-spectrum-sharp-family} 的设定下，若四阶中心矩有限，则对任意
\[
\alpha\in(0,\infty)\setminus\{1\},
\]
记
\[
D_\alpha(\tilde h_t\mid \tilde g_t)
:=\frac{1}{\alpha-1}\log\int_{\RR}\tilde h_t(x)^\alpha \tilde g_t(x)^{1-\alpha}\,\dd x,
\]
有
\[
\lim_{t\to\infty}\frac{t^4}{\sigma^4}\,D_\alpha(\tilde h_t\mid \tilde g_t)=\frac{\alpha}{8}.
\]
\end{conclusion}
\begin{proof}[证明]
令 \(\delta_t:=\tilde h_t/\tilde g_t-1\)。由二阶正规形有 \(\delta_t=O(t^{-2})\) 且
\(\int \tilde g_t\delta_t\,\dd x=0\)。
于是
\[
\int \tilde h_t^\alpha \tilde g_t^{1-\alpha}\,\dd x
=\int \tilde g_t(1+\delta_t)^\alpha\,\dd x
=1+\frac{\alpha(\alpha-1)}{2}\int \tilde g_t\delta_t^2\,\dd x+o(t^{-4}).
\]
再用
\(
\int \tilde g_t\delta_t^2\,\dd x=\sigma^4/(4t^4)+o(t^{-4})
\)
与 \(\log(1+u)=u+o(u)\)，得到
\[
D_\alpha(\tilde h_t\mid \tilde g_t)
=\frac{\alpha}{2}\cdot\frac{\sigma^4}{4t^4}+o(t^{-4})
=\frac{\alpha\sigma^4}{8t^4}+o(t^{-4}),
\]
从而结论成立。证毕。
\end{proof}

\begin{conclusion}[KL 与 \texorpdfstring{$L^1$}{L1} 平方等价的渐近常数]\label{con:xi-poisson-kl-l1-square-equivalence-2pi2-27}
在结论 \ref{con:xi-poisson-lp-spectrum-sharp-family} 的条件下，
\[
\lim_{t\to\infty}
\frac{D_{\mathrm{KL}}(\tilde h_t\mid \tilde g_t)}
{\|\tilde h_t-\tilde g_t\|_{L^1(\RR)}^2}
=\frac{2\pi^2}{27}.
\]
等价地，若
\(
\mathrm{TV}(\tilde h_t,\tilde g_t):=\frac12\|\tilde h_t-\tilde g_t\|_{L^1},
\)
则
\[
D_{\mathrm{KL}}(\tilde h_t\mid \tilde g_t)
\sim \frac{8\pi^2}{27}\,\mathrm{TV}(\tilde h_t,\tilde g_t)^2
\qquad(t\to\infty).
\]
\end{conclusion}
\begin{proof}[证明]
由定理 \ref{thm:xi-cauchy-poisson-tv-sharp-constant} 与
\ref{thm:xi-cauchy-poisson-kl-sharp-asymptotics}，
\[
\|\tilde h_t-\tilde g_t\|_{L^1}
=\frac{3\sqrt3}{4\pi}\frac{\sigma^2}{t^2}+O(t^{-3}),
\qquad
D_{\mathrm{KL}}(\tilde h_t\mid \tilde g_t)=\frac{\sigma^4}{8t^4}+o(t^{-4}),
\]
直接比值化简即得常数 \(\frac{2\pi^2}{27}\)；
TV 版本乘以系数 \(4\) 即得。证毕。
\end{proof}

\begin{conclusion}[反对称剖面分离三阶中心矩]\label{con:xi-poisson-antisymmetric-profile-mu3}
在结论 \ref{con:xi-poisson-lp-spectrum-sharp-family} 的条件下，若四阶中心绝对矩有限，且 \(\mu_3\neq 0\)，记
\[
\mu_3:=\int_{\RR}(\gamma-\bar\gamma)^3\,\dd\tilde\nu(\gamma).
\]
则对任意 \(y\in\RR\) 有一致极限
\[
\lim_{t\to\infty}
\frac{t^4}{\mu_3}
\Bigl(\tilde h_t(\bar\gamma+ty)-\tilde h_t(\bar\gamma-ty)\Bigr)
=\frac{8}{\pi}\frac{y(y^2-1)}{(1+y^2)^4}.
\]
\end{conclusion}
\begin{proof}[证明]
三阶展开给出
\[
\tilde h_t=\tilde g_t+\sigma^2 r_t+\mu_3 s_t+\eta_t,\qquad
s_t:=-\frac16\partial_x^3\mathsf P_t(\cdot-\bar\gamma),
\]
且 \(\|\eta_t\|_\infty=O(t^{-5})\)。
由于 \(\tilde g_t,r_t\) 关于 \(\bar\gamma\) 为偶函数，差分后仅剩奇项：
\[
\tilde h_t(\bar\gamma+ty)-\tilde h_t(\bar\gamma-ty)
=2\mu_3 s_t(\bar\gamma+ty)+O(t^{-5}).
\]
直接微分得到
\[
s_t(\bar\gamma+ty)=\frac{4}{\pi t^4}\frac{y(y^2-1)}{(1+y^2)^4},
\]
故
\[
\frac{t^4}{\mu_3}\Bigl(\tilde h_t(\bar\gamma+ty)-\tilde h_t(\bar\gamma-ty)\Bigr)
=\frac{8}{\pi}\frac{y(y^2-1)}{(1+y^2)^4}+O(t^{-1}),
\]
且余项一致，结论成立。证毕。
\end{proof}

\begin{conclusion}[对称残差分离四阶中心矩]\label{con:xi-poisson-symmetric-residual-mu4}
在结论 \ref{con:xi-poisson-kl-two-term-sharp-explicit} 的条件下，对任意 \(y\in\RR\) 有一致极限
\[
\lim_{t\to\infty}\frac{t^5}{\mu_4}
\Bigl[
\tilde h_t(\bar\gamma+ty)+\tilde h_t(\bar\gamma-ty)
-2\tilde g_t(\bar\gamma+ty)-2\sigma^2 r_t(\bar\gamma+ty)
\Bigr]
=\frac{2}{\pi}\frac{1-10y^2+5y^4}{(1+y^2)^5}.
\]
\end{conclusion}
\begin{proof}[证明]
四阶展开写作
\[
\tilde h_t=\tilde g_t+\sigma^2 r_t+\mu_3 s_t+\mu_4 q_t+\rho_t,\qquad
q_t:=\frac{1}{24}\partial_x^4\mathsf P_t(\cdot-\bar\gamma),
\]
且 \(\|\rho_t\|_\infty=O(t^{-6})\)。
对称化后奇项 \(s_t\) 抵消，于是
\[
\tilde h_t(\bar\gamma+ty)+\tilde h_t(\bar\gamma-ty)-2\tilde g_t(\bar\gamma+ty)-2\sigma^2 r_t(\bar\gamma+ty)
=2\mu_4 q_t(\bar\gamma+ty)+O(t^{-6}).
\]
又
\[
q_t(\bar\gamma+ty)=\frac{1}{\pi t^5}\frac{1-10y^2+5y^4}{(1+y^2)^5},
\]
故乘以 \(t^5/\mu_4\) 即得结论。证毕。
\end{proof}

\begin{conclusion}[同中心 Cauchy 族的信息投影与最优尺度修正]\label{con:xi-poisson-cauchy-information-projection-optimal-scale}
在结论 \ref{con:xi-poisson-lp-spectrum-sharp-family} 的设定下，考虑同中心 Cauchy 族
\[
\tilde g_s(x):=\mathsf P_s(x-\bar\gamma)\qquad(s>0).
\]
取任意尺度 \(s_t\) 满足 \(s_t/t\to1\)，并记
\[
\beta:=\lim_{t\to\infty}t(s_t-t)
\]
（可在任意收敛子列上定义）。令
\[
G(y):=\frac{1}{\pi(1+y^2)},\qquad
u(y):=\frac{3y^2-1}{(1+y^2)^2},\qquad
\phi(y):=\frac{y^2-1}{1+y^2}.
\]
则
\[
D_{\mathrm{KL}}(\tilde h_t\mid \tilde g_{s_t})
=\frac{\sigma^4}{2t^4}
\int_{\RR}G(y)\left(u(y)-\frac{\beta}{\sigma^2}\phi(y)\right)^2\dd y
+o(t^{-4}).
\]
特别地，
\[
\liminf_{t\to\infty}\frac{t^4}{\sigma^4}D_{\mathrm{KL}}(\tilde h_t\mid \tilde g_{s_t})\ge \frac{1}{16},
\]
且当且仅当
\[
s_t=t+\frac{\sigma^2}{2t}+o(t^{-1})
\]
时取等号；此时
\[
D_{\mathrm{KL}}\!\left(\tilde h_t\mid \tilde g_{\,t+\sigma^2/(2t)}\right)
=\frac{\sigma^4}{16t^4}+o(t^{-4}).
\]
\end{conclusion}
\begin{proof}[证明]
该结论与定理 \ref{thm:xi-poisson-optimal-cauchy-rescaling-kl-constant} 同构。
令 \(y=(x-\bar\gamma)/t\)，并写 \(s_t=t+\beta/t+o(t^{-1})\)，则
\[
\frac{\tilde g_{s_t}(\bar\gamma+ty)}{\tilde g_t(\bar\gamma+ty)}
=1+\frac{\beta}{t^2}\phi(y)+O(t^{-4}),
\]
而
\[
\frac{\tilde h_t(\bar\gamma+ty)}{\tilde g_t(\bar\gamma+ty)}
=1+\frac{\sigma^2}{t^2}u(y)+O(t^{-3}).
\]
故
\[
\frac{\tilde h_t}{\tilde g_{s_t}}-1
=\frac{1}{t^2}\left(\sigma^2u-\beta\phi\right)+O(t^{-3}),
\]
代入 KL 二阶展开即可得到主项二次型。
最小化参数 \(c:=\beta/\sigma^2\)：
\[
\Psi(c):=\int G(u-c\phi)^2
=\int Gu^2-2c\int Gu\phi+c^2\int G\phi^2,
\]
并用
\[
\int Gu^2=\frac14,\qquad
\int Gu\phi=\frac14,\qquad
\int G\phi^2=\frac12
\]
得 \(\Psi(c)=\frac14-\frac{c}{2}+\frac{c^2}{2}\)，唯一极小点 \(c=\frac12\)，最小值 \(\Psi_{\min}=\frac18\)。
因此最小 KL 主项为 \(\frac12\cdot\frac18=\frac1{16}\)，对应
\(
s_t=t+\sigma^2/(2t)+o(t^{-1})
\)。
证毕。
\end{proof}

\begin{remark}[同形公开索引排查的范围说明]\label{rem:xi-poisson-profile-divergence-public-index-note}
关于六阶系数组合、最优尺度修正与相关常数的同形公开索引排查，可与注
\ref{rem:xi-poisson-refined-sixth-order-public-index-check}
联读。该类检索结论仅在既定检索语料与参数化口径内成立，不构成形式意义上的排他性证明。
\end{remark}

\endinput

